% abs, norm
\DeclarePairedDelimiter\abs{\lvert}{\rvert}%
\DeclarePairedDelimiter\norm{\lVert}{\rVert}%



%%%%%%%%%%%%%% Theorem boxes 
\tcbuselibrary{theorems}
\tcbset{boxrule=0pt,toprule=1pt}

\newtheoremstyle{theoremdd}% name of the style to be used
  {0.5\baselineskip}% measure of space to leave above the theorem. E.g.: 3pt
  {0.5\baselineskip}% measure of space to leave below the theorem. E.g.: 3pt
  {}% name of font to use in the body of the theorem
  {}% measure of space to indent
  {\bfseries}% name of head font
  {:}% punctuation between head and body
  {.5em}% space after theorem head; " " = normal interword space
  {}
\theoremstyle{theoremdd}

\newcommand{\newcooltheorem}[5]{%
  \makeatletter
  \Crefname{tcb@cnt@#1}{#2}{#2s}%
  %\Crefname{tcb@cnt@#1}{#2}
  \makeatother
  \newtcbtheorem[list inside=#1, number within=chapter]{#1}{#2}{
    % if number within section: add thesection to enumthm cref
    % Example Frame Start
    empty,% Empty previously set parameters
    % Attaching a box requires an overlay
    attach boxed title to top left,
    % Ensures proper line breaking in longer titles
    minipage boxed title,
    % (boxed title style requires an overlay)
    boxed title style={empty,size=minimal, 
    toprule=0pt,top=4pt,left=3mm},
    % title formatting
    coltitle=black,fonttitle=\scshape,description font=\normalfont\bfseries\sffamily, separator sign={. },
    % spacing
    before=\par\medskip\noindent,parbox=false,boxsep=0pt,left=3mm,right=3mm,top=5pt,breakable,pad at break=0mm,
    before upper=\csname @totalleftmargin\endcsname0pt, % Use instead of parbox=true. 
    %This ensures parskip is inherited by box.
    %
    % Handles box when it exists on one page only
    overlay unbroken={\draw[#4,line width=.5pt] ([xshift=-0pt]title.north west) -- ([xshift=-0pt]frame.south west); 
    \begin{scope}[on background layer]
      \fill[#5] 
      (frame.south west) -- 
      (title.north west) -- 
      ([xshift=3mm]title.north east) --
      (frame.south east);
      % \node (box) [fill=blue!20, fit = {(frame.south west),
      % (title.north west),
      % (title.north east),
      % (frame.south east)}{};
    \end{scope}[on background layer]
    },
    % Handles multipage box: first page
    overlay first={\draw[#4,line width=.5pt] ([xshift=-0pt]title.north west) -- ([xshift=-0pt,yshift=-3mm]frame.south west); 
    \begin{scope}[on background layer]
      \fill[#5] 
      ([yshift=-3mm]frame.south west) -- 
      (title.north west) -- 
      ([xshift=3mm]title.north east) --
      ([yshift=-3mm]frame.south east);
      \end{scope}[on background layer]
    },
    % Handles multipage box: middle page
    overlay middle={\draw[#4,line width=.5pt] ([xshift=-0pt, yshift=3mm]frame.north west) -- ([xshift=-0pt]frame.south west); 
    \begin{scope}[on background layer]
      \fill[#5] 
      (frame.south west) -- 
      ([yshift=3mm]frame.north west) -- 
      ([yshift=3mm]frame.north east) --
      (frame.south east);
      \end{scope}[on background layer]
    },
    % Handles multipage box: last page
    overlay last={\draw[#4,line width=.5pt] ([xshift=-0pt, yshift=3mm]frame.north west) -- ([xshift=-0pt]frame.south west); 
    \begin{scope}[on background layer]
      \fill[#5] 
      (frame.south west) -- 
      ([yshift=3mm]frame.north west) -- 
      ([yshift=3mm]frame.north east) --
      (frame.south east);
      \end{scope}[on background layer]
      },%
    }{#3}%
}%

\newcooltheorem{mydefinition}{Definition}{def}{myblue}{white} %{skyblue!70!black}   % todo rot/orange (?)
\newcooltheorem{myexample}{Example}{ex}{mygreen}{mygreen, opacity=0.07} %{bluishgreen!70!black}      % todo grün
\newcooltheorem{mytheorem}{Theorem}{theo}{myred}{white} %{reddishpurple!70!black}  %todo blau
\newcooltheorem{mylemma}{Lemma}{theo}{myred}{white}
\newcooltheorem{myproposition}{Proposition}{theo}{myred}{white}
\newcooltheorem{mycorollary}{Corollary}{theo}{myred}{white} 
\newcooltheorem{myproblem}{Problem}{prob}{mypink}{mypink, opacity=0.1}

\theoremstyle{remark}
\newtheorem{remark}{Remark}[chapter]

% footnotemarks inside align
\makeatletter
\let\original@footnotemark\footnotemark
\newcommand{\align@footnotemark}{%
  \ifmeasuring@
    \chardef\@tempfn=\value{footnote}%
    \original@footnotemark
    \setcounter{footnote}{\@tempfn}%
  \else
    \iffirstchoice@
      \original@footnotemark
    \fi
  \fi}
\pretocmd{\start@align}{\let\footnotemark\align@footnotemark}{}{}
\makeatother

\newcommand{\defeq}{\vcentcolon=}
\newcommand{\eqdef}{=\vcentcolon}
\newcommand{\N}{\mathbb{N}}
\newcommand{\R}{\mathbb{R}}
\newcommand{\Z}{\mathbb{Z}}
\newcommand{\Q}{\mathbb{Q}}
\newcommand{\C}{\mathbb{C}}
\newcommand{\Pred}{\operatorname{Pred}}
\newcommand{\Succ}{\operatorname{Succ}}
\newcommand{\Reach}{\operatorname{Reach}}
\newcommand{\Cone}{\operatorname{Cone}}
\newcommand{\satr}{\operatorname{satr}}
\newcommand{\watr}{\operatorname{watr}}
\newcommand{\ASReach}{\operatorname{ASReach}}
\renewcommand{\Pr}{\operatorname{Pr}}
\newcommand\CC{
    C++
}
\renewcommand{\set}[1]{\mathopen{}\left\{\;\!#1\;\!\right\}\mathclose{}}

\def\until{\mathsf{U}}

\newcommand{\mycomment}[2][]{\todo[size=\footnotesize, #1]{#2}}
\newcommand{\mytodo}[2][]{\todo[color=red, size=\footnotesize, #1]{TODO: #2}}

% multiline stuff for algpseudocode
\iffalse
\newcommand\CONDITION[2]%
   {\begin{tabular}[t]{@{}l@{}}#1#2\end{tabular}}
\algdef{SE}[WHILE]{While}{EndWhile}[1]%
   {\algorithmicwhile\ \CONDITION{#1}{\ \algorithmicdo}}%
   {\algorithmicend\ \algorithmicwhile}
\algdef{SE}[FOR]{For}{EndFor}[1]%
   {\algorithmicfor\ \CONDITION{#1}{\ \algorithmicdo}}%
   {\algorithmicend\ \algorithmicfor}
\algdef{S}[FOR]{ForAll}[1]%
   {\algorithmicforall\ \CONDITION{#1}{\ \algorithmicdo}}
\algdef{SE}[REPEAT]{Repeat}{Until}%
   {\algorithmicrepeat}[1]%
   {\algorithmicuntil\ \CONDITION{#1}{}}
\algdef{SE}[IF]{If}{EndIf}[1]%
   {\algorithmicif\ \CONDITION{#1}{\ \algorithmicthen}}%
   {\algorithmicend\ \algorithmicif}
\algdef{C}[IF]{IF}{ElsIf}[1]%
   {\algorithmicelse\ \algorithmicif\ \CONDITION{#1}{\ \algorithmicthen}}
   
\makeatletter
\ifthenelse{\equal{\ALG@noend}{t}}%
   {\algtext*{EndWhile}%
    \algtext*{EndFor}%
    \algtext*{EndLoop}%
    \algtext*{EndIf}%
   }{}%
\makeatother
\fi
